\lesson{goal}{把想法变成可执行的任务}
\diagram{p_goal}
\topic{先把结果说清楚}
\begin{enumerate}
\item \textbf{说明目标与依据。}写清谁要使用、要完成什么操作，附上已有资料或认可的参考。描述希望达到的结果，不必预先规定所有技术步骤。
\item \textbf{划定第一版。}区分现在必须完成与以后再做的内容。想法模糊时，让 AI 每次问一个关键问题，逐步明确场景、边界和取舍；已明确的要求无需反复确认。
\item \textbf{约定完成标准。}把“好用、完整、专业”改成可以观察的行为：哪些操作应成功、失败时怎样处理、哪些内容需要保留。先检查这一版需求是否遗漏，再安排实现。
\end{enumerate}
\begin{prompt}{需求整理与最小方案}
我要让[使用者]在[场景]完成[目标]。项目和参考资料在[位置；尚无则注明]。本版必须完成[范围]，暂不做[内容]；允许修改[范围]。\par
先核对现状，保留原始要求，在现有需求文档或当前功能的 spec.md 中写清目标、关键操作、完成标准、必须保留的行为和未知项。能查证的自行查证；关键需求不清时逐个询问并给建议，标明未确认的假设。\par
提出最小方案和有检查方法的小任务。本轮只明确需求与初步安排，不改代码；技术选择待调研。沿用已有文档或 Spec Kit 步骤；材料不可读先说明缺口，无法写文件则返回草稿，不声称已保存。
\end{prompt}
\principle{把重要选择留给人，把实现空间留给 AI}{需求规格（Spec）写清目标、行为、约束和验收标准；计划写“准备怎么做”。可以先与 AI 讨论想法，再整理成规格；人重点审核目标、取舍和验收，技术方案经调研后再定。标准明确有利于自动执行，却不能保证理解和实现都正确，结果仍需核验。目标清楚的小改动可直接执行，不必另写长计划。}
